\section{An unbounded obstruction to uniform cubic defect descent}
\label{sec:aa-cubic-amplification}
This section gives an unconditional obstruction to one precise child of the
power-descent route. It concerns actual radicals of actual primitive triples,
and is stronger than an obstruction to an independently estimated lifting
factor. The parent route using the defect of the seed remains open.

For $x\ge2$ put
\[
 S_x=(1,x-1,x),\qquad L_x=(1,x^3-1,x^3),\qquad
 R_1=\operatorname{rad}(x(x-1)),\quad
 R_3=\operatorname{rad}(x(x^3-1)).
\]
For an integer $m\ge1$ write
\[
 \mathcal Q_m(a,b,c)=c^m/\operatorname{rad}(abc)^{m+1}.
\]
Both displayed triples are positive and primitive.

\begin{lemma}[A direct radical compression bound]
\label{lem:aa-compression}
Suppose that $h\ge1$ and $7^h\mid x^2+x+1$. Then
\[
 7^{h-1}R_3\le R_1(x^2+x+1).
\]
If also $x<4\cdot7^h$, then $R_3<84xR_1$.
\end{lemma}
\begin{proof}
Set $G=x^2+x+1$. The factorization $x^3-1=(x-1)G$ gives
$R_3\le R_1\operatorname{rad}(G)$. Write $G=7^v u$ with
$v\ge h$ and $7\nmid u$. Then
\[
 7^{h-1}\operatorname{rad}(G)
 =7^h\operatorname{rad}(u)\le7^h u\le G.
\]
This proves the first inequality without requiring any unproved information
about the other prime factors of $G$. Since $x\ge2$, $G\le3x^2$.
Also $x<4\cdot7^h=28\cdot7^{h-1}$, so
$R_3\le3x^2R_1/7^{h-1}<84xR_1$.
\end{proof}

\begin{lemma}[Infinitely many normalized compressed bases]
\label{lem:aa-bases}
For every integer $h\ge2$ there exists an integer $x_h$ such that
\[
 2\le x_h<4\cdot7^h,\qquad x_h\equiv2\pmod4,\qquad
 x_h\equiv2\pmod7,\qquad 7^h\mid x_h^2+x_h+1.
\]
Every $x_h$ is not a perfect power of exponent at least two, and
$x_h\longrightarrow\infty$ as $h\longrightarrow\infty$.
\end{lemma}
\begin{proof}
The polynomial $F(X)=X^2+X+1$ satisfies $F(2)=7$ and
$F'(2)=5\not\equiv0\pmod7$. A completely elementary digit lift suffices:
if $r$ is a root modulo $7^j$, choose the unique $d\in\{0,\ldots,6\}$
such that
\[
 F(r)/7^j+d(2r+1)\equiv0\pmod7.
\]
Then $r+d7^j$ is a root modulo $7^{j+1}$ and is still $2$ modulo seven.
Induction gives a root $r_h$ modulo $7^h$. The Chinese remainder theorem
gives the unique $x_h\in[0,4\cdot7^h)$ satisfying
$x_h\equiv r_h\pmod{7^h}$ and $x_h\equiv2\pmod4$.
Its latter residue implies $x_h\ge2$ and $v_2(x_h)=1$, excluding every
nontrivial perfect power. Finally $7^h\le F(x_h)$; since $F$ is bounded
on every bounded interval of nonnegative integers, this proves
$x_h\to\infty$. The sequence need not be monotone.
\end{proof}

\begin{theorem}[Unbounded amplification of the actual cleared defect]
\label{thm:aa-amplification}
For every integer $m\ge2$, the normalized bases in
Lemma~\ref{lem:aa-bases} satisfy
\[
 \boxed{\quad
 \frac{\mathcal Q_m(L_{x_h})}{\mathcal Q_m(S_{x_h})}
 >\frac{x_h^{m-1}}{84^{m+1}}\longrightarrow\infty.
 \quad}
\]
Consequently there is no finite constant $A_m$ such that
$\mathcal Q_m(L_x)\le A_m\mathcal Q_m(S_x)$ for every integer
$x\ge2$ which is not a nontrivial perfect power.
\end{theorem}
\begin{proof}
The exact quotient of the two positive defects is
\[
 \frac{\mathcal Q_m(L_x)}{\mathcal Q_m(S_x)}
 =x^{2m}(R_1/R_3)^{m+1}.
\]
Lemma~\ref{lem:aa-compression} bounds this from below by
$x^{2m}/(84x)^{m+1}=x^{m-1}/84^{m+1}$. For $m\ge2$ this tends to
infinity along Lemma~\ref{lem:aa-bases}. All arithmetic premises of the
refuted estimate, including normalization, hold throughout this family.
\end{proof}

\begin{corollary}[The entire relevant real-exponent interval]
\label{cor:aa-real-amplification}
For $\varepsilon>0$ define
$D_\varepsilon(a,b,c)=c/\operatorname{rad}(abc)^{1+\varepsilon}$.
For every fixed $0<\varepsilon<1$,
\[
 \frac{D_\varepsilon(L_{x_h})}{D_\varepsilon(S_{x_h})}
 >\frac{x_h^{1-\varepsilon}}{84^{1+\varepsilon}}
 \longrightarrow\infty.
\]
\end{corollary}
\begin{proof}
The exact quotient is $x_h^2(R_1/R_3)^{1+\varepsilon}$.
Apply Lemma~\ref{lem:aa-compression}. The exponent $1-\varepsilon$ is
positive.
\end{proof}

\paragraph{Scope and the surviving bottleneck.}
The theorem refutes even a bounded-loss universal cubic defect comparison,
not only a universal nonincrease comparison. It does not prove the target
defects themselves unbounded: their seed defects may tend to zero. In
particular the precise remaining obligation in a seed-sensitive approach is
to control the product
\[
 \mathcal Q_m(S_x)\,
 x^{2m}(R_1/R_3)^{m+1}
 =\mathcal Q_m(L_x).
\]
Treating this identity as a new independent estimate would be circular.
The counterfamily excludes only the universal relative-comparison child;
power descent with arithmetic seed credit and all other parent routes remain
active. The endpoint $\varepsilon=1$ is not decided by this lower bound.
Neither this theorem nor finite computations below prove or disprove standard
$abc$.

\paragraph{Formal and computational scope.}
The companion \texttt{CubicAmplification.lean} was checked using Lean 4.32.0
with warnings as errors. Its nine named declarations prove the cofactor cap
and monotonicity, an explicit normalized root above any prescribed bound, the
compression-to-amplification implications, and the impossibility of a uniform
relative bound conditional on the displayed compression inequalities.
The imported 29-declaration power-descent module was freshly compiled with
the same compiler. The reported axiom dependencies are limited to
\texttt{propext}, \texttt{Classical.choice} and \texttt{Quot.sound}.
This is a scoped check; the actual radical function and its prime-factorization
inequality in Lemma~\ref{lem:aa-compression} remain ordinary proofs, and are
not claimed as formalized by this module. An independent exact-integer replay
checks six fully factorized rows and 100 lifting levels. These finite checks
neither certify an asymptotic estimate nor replace the infinite argument.
